% Grants, supervision, committees, awards, service and courses as biblatex
% records (data model: uwhealthcv.dbx).
% Part of the template, loaded by uwhealthcv.sty after aprcv.sty. The CV
% keeps the records in .bib files and prints them with:
%
%   \cvowner{Name}               whose CV this is, set once: \cvname prints
%                                it (NAME/DATE), \me in a grant prints it in
%                                bold, the running header shows it, and the
%                                PDF title and author use it
%   \printgrants                 the RESEARCH GRANTS AND CONTRACTS table,
%                                newest first (by start, then end year):
%                                the grants with no status (awarded)
%   \printgrants[Submitted]      the same table for the grants whose status
%                                is Submitted (any status works: a separate
%                                table under its own heading for applications
%                                awaiting a decision)
%   \printsupervision{PhD}       one numbered list of students under a
%                                degree label: completed by date, then
%                                those in progress, by family name; a
%                                student with started prints the years
%                                (2019 -- 2022) or In progress (since 2023)
%   \printcommitteecounts{MSc, PhD}
%                                thesis committees: one line per year with
%                                the number in each group (see the
%                                committees part below, which also has
%                                \printcommittees and \cvcommitteegroup)
%   \printawards                 ACADEMIC AWARDS AND DISTINCTIONS: newest
%                                first, in file order within a year
%   \printawards[research]       only the awards whose level names it (like
%                                \printservice: level may name several), so
%                                a CV can group its awards under headings
%                                (research, teaching, recognition, trainee
%                                ...) or print a group in another section
%   \printservice{school}        one list of @service records: those whose
%                                level names it (a record may name several),
%                                most recent first (open-ended first, then by
%                                end year, then start), in file order otherwise
%   \printcourses                COURSES TAUGHT: a table of @course records,
%                                newest term first (term, course, level,
%                                enrolment)
%   \printcourses[scp]           the same with two more columns: the course
%                                evaluation scores and their comparison
%                                averages, headed \cvcoursescplabel and
%                                \cvcourseaveragelabel (\renewcommand them);
%                                the averages column only if some course
%                                has scpaverage
%
% In the APR version a grant whose tenure overlaps the review period is
% highlighted as research (the whole table row); a student who completed in the
% period, or is in progress, is highlighted as teaching; an award in the
% period is highlighted as an award; a service role whose years overlap the
% period is highlighted as service; a course taught in the period is
% highlighted as teaching (the whole table row).

\RequirePackage{siunitx}

\def\cv@owner{}
\newcommand\cvowner[1]{\def\cv@owner{#1}%
  \hypersetup{pdfauthor={#1}, pdftitle={Curriculum Vitae: #1}}}% PDF/XMP metadata
% \cvname: the name set with \cvowner, e.g. \cvsection{NAME/DATE:} \cvname, ...
\newcommand\cvname{\cv@owner}

\DeclareFieldFormat{cvplain}{#1}
\DeclareFieldFormat{cvcosupervisor}{\mkbibparens{co-supervised with #1}}
\DeclareFieldFormat{cvservicenote}{\mkbibemph{#1}}

%% ----------------------------------------------------------------- grants
% The grants print as a longtable, like a hand-made table: headings repeat
% on each page, and a grant never splits across pages. biblatex cannot print
% entries into table cells, so the grant driver collects one table row per
% grant into \cv@rows (printing nothing), and \printgrants then sets them.
%
% Column widths: names, agency and short title wrap; amount and tenure don't.
% The table is exactly \linewidth wide and lines up with the text margins: no
% padding at the outer edges, 2\cvgrantcolsep between columns. Amount and
% tenure take their widest entry, and the short title takes what is left
% (\printgrants works it out), so a wide amount narrows the title column.
\newlength\cvgrantnamewidth   \setlength\cvgrantnamewidth{4.15cm}
\newlength\cvgrantagencywidth \setlength\cvgrantagencywidth{4.15cm}
\newlength\cvgranttitlewidth
\newlength\cvgrantcolsep      \setlength\cvgrantcolsep{4pt}

% every investigator, never truncated to "et al." (biblatex's default
% truncates lists longer than three)
\ExecuteBibliographyOptions[grant]{maxitems=999,minitems=999}

% collect the investigators, one per line, into \cv@inv
\DeclareListFormat{cvcollect}{%
  \ifnumgreater{\value{listcount}}{1}{\gappto\cv@inv{\newline}}{}%
  \gappto\cv@inv{#1}}

% \cv@field{name}: a field's text, unexpanded (empty when it is not set)
\newcommand\cv@field[1]{%
  \ifcsundef{abx@field@#1}{}%
    {\unexpanded\expandafter\expandafter\expandafter{%
       \csname abx@field@#1\endcsname}}}

\gdef\cv@rows{}
\DeclareBibliographyDriver{grant}{%
  \gdef\cv@inv{}%
  \printlist[cvcollect]{investigators}%
  \aprsetentrycolour
  \ifdefempty\cv@rows{}{\gappto\cv@rows{\\}}% blank row between grants
  \xappto\cv@rows{%
    \noexpand\cv@grantline
      {\expandonce\aprentrycolour}%
      {\cv@field{role}}%
      {\expandonce\cv@inv}%
      {\cv@field{agency}}%
      {\cv@field{amount}}%
      {\cv@field{year}}%
      {\cv@field{endyear}}%
      {\cv@field{title}}%
      {\expandonce\aprentrymark}}}

% \cv@grantline{colour}{role}{names}{agency}{amount}{year}{endyear}{title}{mark}
% \me is the CV's owner, in bold, with the role; mark is the APR dots
% style's marker (\aprentrymark)
\newcommand\cv@grantline[9]{%
  \ifx\relax#1\relax\else\rowcolor{#1}\fi
  #9\def\me{\textbf{\cv@owner\ifx\relax#2\relax\else\space(#2)\fi}}#3 &
  #4 &
  \cv@amount{#5} &
  \cv@tenure{#6}{#7} &
  #8 \\*}
\newcommand\cv@amount[1]{\$\num[group-separator={,},group-minimum-digits=4]{#1}}
\newcommand\cv@tenure[2]{#1\ifx\relax#2\relax\else\ -- #2\fi}

% \cv@widest\len{text}: widen \len (globally) to fit text
\newcommand\cv@widest[2]{%
  \settowidth\@tempdima{#2}%
  \ifdim\@tempdima>#1\global#1\@tempdima\fi}
\newlength\cv@amountwd
\newlength\cv@tenurewd
% measure the amount and tenure columns over the collected rows
\newcommand\cv@measureline[9]{%
  \cv@widest\cv@amountwd{\cv@amount{#5}}%
  \cv@widest\cv@tenurewd{\cv@tenure{#6}{#7}}}

\DeclareSortingTemplate{cvgrants}{
  \sort[direction=descending]{\field{year}}
  \sort[direction=descending]{\field{endyear}\field{year}}
  \sort{\field{title}}
  \sort{\field{agency}}
}

\defbibenvironment{cvcollect}{}{}{}

% \printgrants prints the grants with no status; \printgrants[Submitted] those
% whose status is Submitted (applications awaiting a decision)
\def\cv@grantstatus{}
\defbibcheck{cvgrantstatus}{%
  \ifdefempty\cv@grantstatus
    {\iffieldundef{status}{}{\skipentry}}%
    {\iffieldequals{status}{\cv@grantstatus}{}{\skipentry}}}

\NewDocumentCommand\printgrants{o}{%
  % end a paragraph a \cvsubheading left open before collecting the rows, so
  % the tagged PDF closes the heading there (the table is not part of it)
  \ifvmode\else\par\fi
  \IfValueTF{#1}{\def\cv@grantstatus{#1}}{\def\cv@grantstatus{}}%
  \gdef\cv@rows{}%
  \setbox\z@\vbox{% collect the rows; nothing is printed here
    \begin{refcontext}[sorting=cvgrants]%
      \printbibliography[type=grant,check=cvgrantstatus,env=cvcollect,heading=none]%
    \end{refcontext}}%
  \begingroup
  \setlength\tabcolsep{\cvgrantcolsep}%
  \global\cv@amountwd\z@
  \global\cv@tenurewd\z@
  \cv@widest\cv@amountwd{\underline{\$ Amount}}%
  \cv@widest\cv@tenurewd{\underline{Tenure}}%
  {\def\\{}\let\cv@grantline\cv@measureline\cv@rows}%
  \setlength\cvgranttitlewidth{\dimexpr\linewidth-\cvgrantnamewidth
    -\cvgrantagencywidth-\cv@amountwd-\cv@tenurewd-8\tabcolsep\relax}%
  \cv@block\setlength\LTpre{\z@}% one blank line after the heading (uwhealthcv.sty)
  \begin{longtable}[c]{@{}>{\raggedright\arraybackslash}p{\cvgrantnamewidth}
                       >{\raggedright\arraybackslash}p{\cvgrantagencywidth}
                       l l
                       >{\raggedright\arraybackslash}p{\cvgranttitlewidth}@{}}
  \underline{Names of Researcher(s)} & \underline{Agency} & \underline{\$ Amount} &
    \underline{Tenure} & \underline{Short Title} \\
  \endhead
  \cv@rows
  \end{longtable}%
  \endgroup}

%% ------------------------------------------------------------ supervision
\aprtypecategory{supervision}{teaching}

% full given names, whatever giveninits says for publications
\DeclareNameFormat{cvfull}{%
  \nameparts{#1}%
  \ifdefvoid\namepartgiven{}{\namepartgiven\space}%
  \ifdefvoid\namepartprefix{}{\namepartprefix\space}%
  \namepartfamily
  \ifdefvoid\namepartsuffix{}{\space\namepartsuffix}}

\newcommand\cv@monthname[1]{%
  \ifcase#1\or January\or February\or March\or April\or May\or June\or
  July\or August\or September\or October\or November\or December\fi}

% Name (Program), date (or status), Title. Note (co-supervised with X) [Position]
% With started (the year the student began): 2019 -- 2022, or
% In progress (since 2023).
\DeclareBibliographyDriver{supervision}{%
  \usebibmacro{begentry}%
  \printnames[cvfull]{author}%
  \setunit{\space}%
  \printfield[parens]{program}%
  \setunit{\addcomma\space}%
  \iffieldint{year}
    {\printtext{%
       \iffieldundef{started}{}{\thefield{started}\ -- }%
       \iffieldint{month}{\cv@monthname{\thefield{month}}\space}{}%
       \thefield{year}}}
    {\printfield[cvplain]{status}%
     \iffieldundef{started}{}{\printtext{\space(since \thefield{started})}}}%
  \setunit{\addcomma\space}%
  \printfield[cvplain]{title}%
  \setunit{\addperiod\space}%
  \printfield[cvplain]{note}%
  \setunit{\space}%
  \printfield[cvcosupervisor]{cosupervisor}%
  \setunit{\space}%
  \printfield[brackets]{position}%
  \usebibmacro{finentry}}

\DeclareSortingTemplate{cvsupervision}{
  \sort{\field{year}\literal{9999}}
  \sort{\field[padside=left,padwidth=2,padchar=0]{month}\literal{99}}
  \sort{\field{author}}
}

% The list starts right after its degree label, so LaTeX treats the first
% \item specially and then clears \if@noparitem. biblatex typesets each
% entry in a group, which would undo that and treat every item as the
% first (the labels pile up), so clear it globally.
\newcommand\cv@degreeitem{\item\global\let\if@noparitem\iffalse}
\defbibenvironment{cvdegree}
  {\begin{enumerate}[label=\arabic*.]%
   \renewcommand*{\finentrypunct}{}}
  {\end{enumerate}}
  {\cv@degreeitem}

\def\cv@degree{}
\defbibcheck{cvdegree}{\iffieldequals{degree}{\cv@degree}{}{\skipentry}}

\newcommand\printsupervision[1]{%
  \cv@degreelist% the degree's label column (uwhealthcv.sty)
    \item[#1]%
    \def\cv@degree{#1}%
    \begin{refcontext}[sorting=cvsupervision]%
      \printbibliography[type=supervision,check=cvdegree,env=cvdegree,heading=none]%
    \end{refcontext}%
  \end{enumerate}}

%% ----------------------------------------------------------------- awards
\aprtypecategory{award}{award}

% [count x] Title
\DeclareBibliographyDriver{award}{%
  \usebibmacro{begentry}%
  \iffieldint{count}
    {\printtext{\thefield{count} $\times$}\setunit{\space}}
    {}%
  \printfield[cvplain]{title}%
  \usebibmacro{finentry}}

% sorted with cvservice (below): newest year first, then file order

% the same layout as cvlist (uwhealthcv.sty)
\defbibenvironment{cvawards}
  {\cv@block\cv@setcm\cv@nobullets\begin{itemize}[leftmargin=\cv@cm, parsep=5pt]%
   \setlength{\labelsep}{0em}\setlength{\itemsep}{-5pt}%
   \renewcommand*{\finentrypunct}{}}
  {\end{itemize}}
  {\item[]}

% \printawards[level]: only the awards whose level names it (the cvlevel
% check, below, that \printservice uses)
\NewDocumentCommand\printawards{o}{%
  \begin{refcontext}[sorting=cvservice]%
    \IfValueTF{#1}
      {\def\cv@level{#1}%
       \printbibliography[type=award,check=cvlevel,env=cvawards,heading=none]}%
      {\printbibliography[type=award,env=cvawards,heading=none]}%
  \end{refcontext}}

%% ---------------------------------------------------------------- service
\aprtypecategory{service}{service}

% Title Years Role: the years field as written, else the date as
% 2023 -- 2025 (an open-ended date, 2021/, prints 2021 -- Present); the role,
% if any, after them; the note, if any, on the next line in italics
\DeclareBibliographyDriver{service}{%
  \usebibmacro{begentry}%
  \printfield[cvplain]{title}%
  \setunit{\space}%
  \iffieldundef{years}
    {\iffieldint{year}
       {\printtext{%
          \thefield{year}%
          \iffieldint{endyear}
            {\iffieldequalstr{endyear}{}
               {\ -- Present}%
               {\iffieldsequal{year}{endyear}{}{\ -- \thefield{endyear}}}}
            {}}}
       {}}
    {\printfield[cvplain]{years}}%
  \setunit{\space}%
  \printfield[cvplain]{role}%
  \iffieldundef{note}{}{\aprlinemark}%
  \setunit{\newline}%
  \printfield[cvservicenote]{note}%
  \usebibmacro{finentry}}

% An open-ended date (2021/) has an empty end year, which biber would sort as
% if the role had ended in its start year; flag those so they sort first.
\DeclareSourcemap{
  \maps[datatype=bibtex]{
    \map{
      \pertype{service}
      \pertype{committee}
      \step[fieldsource=date, match=\regexp{/\s*\Z}, final]
      \step[fieldset=userf, fieldvalue={open}]
    }
  }
}

% \printawards and \printcourses use this sorting too: awards have only a
% year, courses also a month (newest term first). biber writes a full copy
% of every entry into the .bbl for each sorting in use, so each new sorting
% template costs TeX a lot of memory; reuse one where the order allows.
\DeclareSortingTemplate{cvservice}{
  \sort[direction=descending]{\field{userf}}
  \sort[direction=descending]{\field{endyear}\field{year}}
  \sort[direction=descending]{\field{year}}
  \sort[direction=descending]{\field[padside=left,padwidth=2,padchar=0]{month}}
  \sort{\citeorder}
}

% level may name several lists, comma-separated: {discipline, organizing}
\def\cv@level{}
\newcommand\cv@checklevel[1]{\ifdefstring\cv@level{#1}{\def\cv@inlevel{x}\listbreak}{}}
\defbibcheck{cvlevel}{%
  \def\cv@inlevel{}%
  \iffieldundef{level}{}{%
    \edef\cv@levels{\thefield{level}}%
    \expandafter\forcsvlist\expandafter\cv@checklevel\expandafter{\cv@levels}}%
  \ifdefempty\cv@inlevel{\skipentry}{}}

\newcommand\printservice[1]{%
  \def\cv@level{#1}%
  \begin{refcontext}[sorting=cvservice]%
    \printbibliography[type=service,check=cvlevel,env=cvawards,heading=none]%
  \end{refcontext}}

%% ------------------------------------------------------------- committees
% @committee records: thesis and examination committees for students the
% owner doesn't supervise. One record per student (a date range for the years
% served), or a count-only record (count, no author) for a year whose
% committees aren't listed by name.
%
%   \cvcommitteegroup{MSc}{MSc, MASc, MMath}
%                                committees for these degrees are counted and
%                                listed under MSc (a degree in no group is a
%                                group of its own)
%   \printcommitteecounts{MSc, PhD}
%                                one line per year, newest first, with the
%                                number of committees served on that year in
%                                each group, in the order given: "2026 --- 4
%                                MSc, 6 PhD". A committee counts in every year
%                                of its date range (an open range, 2023/, runs
%                                to the current year). Count-only records for
%                                a year and group give that year's number
%                                instead (their counts are added up), so years
%                                with no names can sit beside years with them.
%                                A group with none that year is left out.
%   \printcommittees{PhD}        the named committees in one group, as a
%   \printcommittees[period]{PhD} bulleted list, most recent first (ongoing,
%                                then by end year); [period]: only those whose
%                                years overlap the review period
%
% In the APR version a year line in the review period, and a named committee
% whose years overlap it, are highlighted as teaching. Add a record's key to
% apralsoservice for a second colour (e.g. chairing an examination).
\aprtypecategory{committee}{teaching}

\newcommand\cvcommitteegroup[2]{\forcsvlist{\cv@setcommitteegroup{#1}}{#2}}
\newcommand\cv@setcommitteegroup[2]{\csdef{cv@cgroup@#2}{#1}}
% \cv@groupof{degree}: the degree's group (expandable)
\newcommand\cv@groupof[1]{\ifcsdef{cv@cgroup@#1}{\csuse{cv@cgroup@#1}}{#1}}

\DeclareFieldFormat{cvcommitteetitle}{\mkbibemph{#1}}

% Name (Degree, Program; role, note). Title
\newif\ifcv@counting
\DeclareBibliographyDriver{committee}{%
  \ifcv@counting
    \cv@countentry
  \else
    \usebibmacro{begentry}%
    \printnames[cvfull]{author}%
    \setunit{\space}%
    \printtext[parens]{%
      \printfield[cvplain]{degree}%
      \setunit{\addcomma\space}%
      \printfield[cvplain]{program}%
      \setunit{\addsemicolon\space}%
      \printfield[cvplain]{role}%
      \iffieldundef{role}{}{\setunit{\addcomma\space}}%
      \printfield[cvplain]{note}}%
    \setunit{\addperiod\space}%
    \printfield[cvcommitteetitle]{title}%
    \usebibmacro{finentry}%
  \fi}

% \cv@yearsof: sets \cv@from and \cv@until to the entry's first and last
% year (an open range ends in the current year)
\newcount\cv@from \newcount\cv@until
\newcommand\cv@yearsof{%
  \cv@from=\thefield{year}\relax
  \cv@until=\cv@from
  \iffieldundef{endyear}{}{%
    \iffieldequalstr{endyear}{}%
      {\cv@until=\year}%
      {\cv@until=\thefield{endyear}\relax}}}

% counts are kept per call of \printcommitteecounts, in
% \cv@cn@<call>@<group>@<year> (named) and \cv@co@<call>@<group>@<year>
% (count-only records)
\newcount\cv@ccall
\newcount\cv@cmin \newcount\cv@cmax
\newcommand\cv@cadd[2]{%
  \ifcsdef{#1}%
    {\csxdef{#1}{\the\numexpr\csuse{#1}+#2\relax}}%
    {\csxdef{#1}{\the\numexpr#2\relax}}}
\newcommand\cv@countentry{%
  \iffieldint{year}{%
    \edef\cv@g{\cv@groupof{\thefield{degree}}}%
    \cv@yearsof
    \iffieldint{count}{\cv@until=\cv@from}{}%
    % the range of years to print, before \cv@countyears moves \cv@from on
    \ifnum\cv@from<\cv@cmin \global\cv@cmin=\cv@from \fi
    \ifnum\cv@until>\cv@cmax \global\cv@cmax=\cv@until \fi
    \iffieldint{count}%
      {\cv@cadd{cv@co@\the\cv@ccall @\cv@g @\the\cv@from}{\thefield{count}}}%
      {\cv@countyears}}{}}
\def\cv@countyears{%
  \ifnum\cv@from>\cv@until\else
    \cv@cadd{cv@cn@\the\cv@ccall @\cv@g @\the\cv@from}{1}%
    \advance\cv@from\@ne
    \expandafter\cv@countyears
  \fi}

\newcount\cv@cy
% \cv@countgroup{group}: add "N group" to \cv@line for year \cv@cy
\newcommand\cv@countgroup[1]{%
  \edef\cv@tmp{\the\cv@ccall @#1@\the\cv@cy}%
  \ifcsdef{cv@co@\cv@tmp}%
    {\edef\cv@n{\csuse{cv@co@\cv@tmp}}}%
    {\ifcsdef{cv@cn@\cv@tmp}{\edef\cv@n{\csuse{cv@cn@\cv@tmp}}}{\def\cv@n{0}}}%
  \ifnum\cv@n>\z@
    \ifx\cv@line\@empty\else\appto\cv@line{, }\fi
    \eappto\cv@line{\cv@n\space#1}%
  \fi}

\newcommand\printcommitteecounts[1]{%
  \begingroup
  \global\advance\cv@ccall\@ne
  \global\cv@cmin=\numexpr\year+1\relax
  \global\cv@cmax=\z@
  \cv@countingtrue
  \setbox\z@\vbox{% count; nothing is printed here, so nothing is tagged
    \ifdefined\SuspendTagging\SuspendTagging{committee counts}\fi
    \printbibliography[type=committee,env=cvcollect,heading=none]}%
  \cv@countingfalse
  \ifnum\cv@cmax>\z@
    \begin{cvlist}%
    \cv@cy=\cv@cmax
    \loop\unless\ifnum\cv@cy<\cv@cmin
      \def\cv@line{}%
      \forcsvlist\cv@countgroup{#1}%
      \ifx\cv@line\@empty\else
        \edef\cv@tmp{\noexpand\cvitem\noexpand\hlteaching[\the\cv@cy]%
          {\the\cv@cy\space---\space\expandonce\cv@line}}%
        \cv@tmp
      \fi
      \advance\cv@cy\m@ne
    \repeat
    \end{cvlist}%
  \fi
  \endgroup}

% named committees in one group (\cv@cgroup), and with [period] only those
% whose years overlap the review period
\def\cv@cgroup{}
\newif\ifcv@cperiod
\defbibcheck{cvcommittee}{%
  \ifnameundef{author}{\skipentry}{%
    \edef\cv@g{\cv@groupof{\thefield{degree}}}%
    \ifx\cv@g\cv@cgroup
      \ifcv@cperiod
        \iffieldint{year}{%
          \cv@yearsof
          \edef\cv@tmp{\the\cv@from-\the\cv@until}%
          \expandafter\apr@ifyears\expandafter{\cv@tmp}{}{\skipentry}}{}%
      \fi
    \else
      \skipentry
    \fi}}

\defbibenvironment{cvcommittees}
  {\cv@block\cv@setcm
   \begin{itemize}[label=\textbullet, leftmargin=\cv@cm, labelsep=0.5em, itemsep=0pt]\raggedright
   \renewcommand*{\finentrypunct}{}}
  {\end{itemize}}
  {\item}

\NewDocumentCommand\printcommittees{om}{%
  \begingroup
  \edef\cv@cgroup{#2}%
  \IfValueTF{#1}{\cv@cperiodtrue}{\cv@cperiodfalse}%
  \begin{refcontext}[sorting=cvservice]%
    \printbibliography[type=committee,check=cvcommittee,env=cvcommittees,heading=none]%
  \end{refcontext}%
  \endgroup}

%% ---------------------------------------------------------------- courses
% Like the grants: the driver collects one table row per course into
% \cv@courserows and \printcourses sets them as a longtable. The course
% cell holds the number, title, level and note. Term, enrolment and the score columns take their natural
% width and the course column takes what is left of \linewidth.
\aprtypecategory{course}{teaching}

\newcommand\cvcoursescplabel{Evaluation}
\newcommand\cvcourseaveragelabel{Average}
\newlength\cvcoursewidth

% the term from the month it starts: 1-4 Winter, 5-8 Spring, 9-12 Fall
\newcommand\cv@term[2]{%
  \ifx\relax#2\relax\else
    \ifnum#2<5 Winter\else\ifnum#2<9 Spring\else Fall\fi\fi\space
  \fi#1}

% scores: {Implementation=4.5, Design=4.6} prints one "Item score/5" per line
\newcommand\cv@scores[1]{%
  \begingroup\def\cv@scoresep{}%
  \forcsvlist\cv@scoreitem{#1}%
  \endgroup}
\newcommand\cv@scoreitem[1]{\cv@scoresep\def\cv@scoresep{\newline}%
  \cv@scoresplit#1==\@nil}
\def\cv@scoresplit#1=#2=#3\@nil{#1\ifx\relax#2\relax\else\space#2/5\fi}

% the course cell: number and title (optional), then the level and the note
% in italics, each on its own line
\newcommand\cv@coursecell[4]{#1\ifx\relax#2\relax\else\space#2\fi
  \ifx\relax#3\relax\else\newline\textit{#3}\fi
  \ifx\relax#4\relax\else\newline\textit{#4}\fi}

\gdef\cv@courserows{}
\DeclareBibliographyDriver{course}{%
  \aprsetentrycolour
  \xappto\cv@courserows{%
    \noexpand\cv@courseline
      {\expandonce\aprentrycolour}%
      {\noexpand\cv@term{\cv@field{year}}{\cv@field{month}}}%
      {\noexpand\cv@coursecell{\cv@field{course}}{\cv@field{title}}{\cv@field{level}}{\cv@field{note}}}%
      {\cv@field{enrolment}}%
      {\noexpand\cv@scores{\cv@field{scp}}}%
      {\noexpand\cv@scores{\cv@field{scpaverage}}}%
      {\expandonce\aprentrymark}}}

% \cv@courseline{colour}{term}{course}{enrolment}{scp}{average}{mark}
% Rows end in \\, so the table can break between courses but not inside one.
% mark is the APR dots style's marker (\aprentrymark).
\newcommand\cv@courseline[7]{%
  \ifx\relax#1\relax\else\rowcolor{#1}\fi
  #7#2 & #3 & #4\cv@coursescores{#5}{#6}\\[0.4\baselineskip]}
\newcommand\cv@coursescores[2]{}
\newcommand\cv@coursescores@on[2]{& #1 & #2}
\newcommand\cv@coursescores@scp[2]{& #1}

% measure the natural-width columns over the collected rows
\newlength\cv@termwd
\newlength\cv@enrolwd
\newlength\cv@scpwd
\newlength\cv@avgwd
% \cv@widestscore\len{item}: widen \len to fit one "Item score/5" line
\newcommand\cv@widestscore[2]{\cv@widest#1{\cv@scoresplit#2==\@nil}}
\newcommand\cv@measurecourse[7]{%
  \cv@widest\cv@termwd{#2}%
  \cv@widest\cv@enrolwd{#4}%
  \cv@measurescores\cv@scpwd#5%
  \cv@measurescores\cv@avgwd#6}
% in the rows, #5 and #6 are \cv@scores{list}: measure each item of the list
\def\cv@measurescores#1\cv@scores#2{\forcsvlist{\cv@widestscore#1}{#2}}

\NewDocumentCommand\printcourses{o}{%
  \gdef\cv@courserows{}%
  \setbox\z@\vbox{% collect the rows; nothing is printed here
    \begin{refcontext}[sorting=cvservice]% newest term first
      \printbibliography[type=course,env=cvcollect,heading=none]%
    \end{refcontext}}%
  \begingroup
  \setlength\tabcolsep{\cvgrantcolsep}%
  \global\cv@termwd\z@ \global\cv@enrolwd\z@
  \global\cv@scpwd\z@  \global\cv@avgwd\z@
  \cv@widest\cv@termwd{\underline{Term}}%
  \cv@widest\cv@enrolwd{\underline{Enrolment}}%
  {\let\cv@courseline\cv@measurecourse\cv@courserows}%
  \IfValueTF{#1}%
    {\cv@widest\cv@scpwd{\underline{\cvcoursescplabel}}%
     \ifdim\cv@avgwd=\z@ % no course has averages: leave out their column
       \let\cv@coursescores\cv@coursescores@scp
       \setlength\cvcoursewidth{\dimexpr\linewidth-\cv@termwd-\cv@enrolwd
         -\cv@scpwd-6\tabcolsep\relax}%
       \def\cv@colspec{@{}l>{\raggedright\arraybackslash}p{\cvcoursewidth}r
         >{\raggedright\arraybackslash}p{\cv@scpwd}@{}}%
       \def\cv@head{\underline{Term} & \underline{Course} & \underline{Enrolment} &
         \underline{\cvcoursescplabel}}%
     \else
       \cv@widest\cv@avgwd{\underline{\cvcourseaveragelabel}}%
       \let\cv@coursescores\cv@coursescores@on
       \setlength\cvcoursewidth{\dimexpr\linewidth-\cv@termwd-\cv@enrolwd
         -\cv@scpwd-\cv@avgwd-8\tabcolsep\relax}%
       \def\cv@colspec{@{}l>{\raggedright\arraybackslash}p{\cvcoursewidth}r
         >{\raggedright\arraybackslash}p{\cv@scpwd}
         >{\raggedright\arraybackslash}p{\cv@avgwd}@{}}%
       \def\cv@head{\underline{Term} & \underline{Course} & \underline{Enrolment} &
         \underline{\cvcoursescplabel} & \underline{\cvcourseaveragelabel}}%
     \fi}%
    {\setlength\cvcoursewidth{\dimexpr\linewidth-\cv@termwd-\cv@enrolwd
       -4\tabcolsep\relax}%
     \def\cv@colspec{@{}l>{\raggedright\arraybackslash}p{\cvcoursewidth}r@{}}%
     \def\cv@head{\underline{Term} & \underline{Course} & \underline{Enrolment}}}%
  \expandafter\cv@coursetable\expandafter{\cv@colspec}%
  \endgroup}
\newcommand\cv@coursetable[1]{%
  \cv@block\setlength\LTpre{\z@}% one blank line after the heading (uwhealthcv.sty)
  \begin{longtable}[c]{#1}
  \cv@head \\
  \endhead
  \cv@courserows
  \end{longtable}}
